% !PRINTED 571
\weberpage{571}
%<<CONT
[\S.\,169, (6)] dadurch eine andere Gestalt erhält. Wir bekommen
also dieselbe Function $\psi_\lambda(\alpha)$, wenn wir $\lambda,\mu$ durch $\lambda f, f$ oder
durch $f, \lambda f$ ersetzen, also [\S.\,169, (8)]

% [resolved 2026-08-06; confirmed 2026-08-07 at 600 dpi] CORRECTED THROUGHOUT THIS LEAF.
% All four sigmas in the two displays carry "t" ABOVE only; what sits beneath each is a short
% horizontal RULE under the sigma, not a second "t". The original transcription "\sum_t^t"
% duplicated the index. Reproduced here as "t" above, with the printed rule recorded in this
% note rather than faked in LaTeX.
\[
\psi_\lambda(\alpha) = \sum^{t} \alpha^{\operatorname{ind} t - (\lambda+1)\operatorname{ind}(t+1)}
= \sum^{t} \alpha^{\lambda \operatorname{ind} t - (\lambda+1)\operatorname{ind}(t+1)}.
\]

Hängen also zwei Zahlen $\lambda,\lambda'$ durch die Congruenz

\weq{1}{
\lambda\lambda' \equiv 1 \pmod{e}
}

mit einander zusammen, so ergiebt sich

\[
\psi_\lambda(\alpha^{\lambda'}) = \sum^{t} \alpha^{\lambda\lambda' \operatorname{ind} t - (\lambda\lambda'+\lambda')\operatorname{ind}(t+1)}
= \sum^{t} \alpha^{\operatorname{ind} t - (\lambda'+1)\operatorname{ind}(t+1)},
\]

also die Formel

\weq{2}{
\psi_\lambda(\alpha^{\lambda'}) = \psi_{\lambda'}(\alpha).
}

Wenn wir sodann in \S.\,169, (5) für $\lambda$ setzen $-\lambda-\mu$ und
also $-\lambda$ für $\lambda+\mu$ (wobei dann vorausgesetzt sein muss, dass
$\lambda$ durch $m$ nicht theilbar ist), so folgt

\[
\frac{(\varepsilon^{-\lambda-\mu}, r)(\varepsilon^\mu, r)}{(\varepsilon^{-\lambda}, r)} = \psi_{-\lambda-\mu,\mu}(\varepsilon).
\]

Die linke Seite multipliciren wir mit

\[
\frac{(\varepsilon^{\lambda+\mu}, r)(\varepsilon^\lambda, r)}{(\varepsilon^{\lambda+\mu}, r)(\varepsilon^\lambda, r)} = 1,
\]

und wenden im Zähler und Nenner die Formel

\[
(\varepsilon^\lambda, r)(\varepsilon^{-\lambda}, r) = (-1)^\lambda n
\]

[\S.\,169, (3)] an, dann folgt

\[
(-1)^\mu \frac{(\varepsilon^\lambda, r)(\varepsilon^\mu, r)}{(\varepsilon^{\lambda+\mu}, r)} = \psi_{-\lambda-\mu,\mu}(\varepsilon),
\]

oder

\[
\psi_{-\lambda-\mu,\mu}(\varepsilon) = (-1)^\mu \psi_{\lambda,\mu}(\varepsilon),
\]

und wenn wir $\lambda,\mu$ durch $\lambda f, f$ ersetzen und $\lambda''$ aus der
Congruenz

\weq{3}{
\lambda + \lambda'' + 1 \equiv 0 \pmod{e}
}

bestimmen,

\weq{4}{
\psi_\lambda(\alpha) = (-1)^f \psi_{\lambda''}(\alpha).
}

Diese Formel kann dann wieder mit der Formel (2) verbunden
werden, und so die Anzahl der für ein bestimmtes $e$ zu berechnenden
Functionen $\psi_\lambda$, deren Anzahl $e-2$ beträgt, etwa auf den
sechsten Theil zurückgeführt werden.

Eine dritte Formel ergiebt sich, wenn man in \S.\,169, (5) für
$\varepsilon$ setzt $\varepsilon^{-1}$ und dann multiplicirt, also
%CONT>>
